%% kapa: exam=1 round=36 year=2025 period=2 form=A subject=appraisal_law subject_ko=감정평가관계법규 seq=16 number=16
%% source: Q-Net 감정평가사 기출문제 (cst00309) artlSeq=5231525 file=raw/1008/5231525/2025년 제36회 감정평가사 1차시험 2교시 문제지 원본.pdf
%% conversion: kordoc 4.24.51 -> tools/gen1.py (auto); review=proofread-sonnet; flags=none
\begin{question}{16}
부동산 가격공시에 관한 법령상 주택가격의 공시에 관한 설명으로 옳지 않은 것은?
\begin{choices}
  \item 군수는 군부동산가격공시위원회의 심의를 거쳐 매년 표준주택가격의 공시기준일 현재 관할 구역 안의 개별주택의 가격을 결정ㆍ공시하여야 한다.
  \item 표준주택으로 선정된 단독주택에 대하여는 개별주택가격을 결정ㆍ공시하지 아니할 수 있고, 해당 주택의 표준주택가격을 개별주택가격으로 본다.
  \item 개별주택가격의 공시사항에는 개별주택의 사용승인일이 포함되어야 한다.
  \item 개별주택가격 및 공동주택가격은 지방자치단체가 지방세 부과 업무와 관련하여 주택의 가격을 산정하는 경우에 그 기준으로 활용될 수 있다.
  \item 표준주택 선정 및 관리에 필요한 세부기준은 중앙부동산가격공시위원회의 심의를 거쳐 국토교통부장관이 정한다.
\end{choices}
\end{question}
\answer{3}
